
\documentclass[11pt]{article}
\usepackage[margin=1in]{geometry}
\usepackage{amsmath,amssymb,booktabs,hyperref,graphicx,parskip,enumitem}
\usepackage{xcolor}
\hypersetup{colorlinks=true,linkcolor=blue!50!black,urlcolor=blue!50!black}
\setlist{nosep}

\title{Replication Report: Simple Coplanar Waveguide Resonator Mask Targeting Multiplexed Superconducting-Qubit Readout\\[2pt]\large OSTI 1983793 \textbar{} Coverage 8/10, Agreement 8/10}
\author{Rick Stevens \and Ollie (AI Assistant)\\\small Argonne National Laboratory}
\date{2026-04-27}

\begin{document}\maketitle

\section{Source paper}
\begin{itemize}
\item \textbf{Title:} Simple Coplanar Waveguide Resonator Mask Targeting Multiplexed Superconducting-Qubit Readout
\item \textbf{Authors:} Yan et al.
\item \textbf{Year:} 2022
\item \textbf{Domain:} Quantum Devices / Superconducting circuits
\item \textbf{Identifier:} OSTI 1983793
\end{itemize}


\section{Replication objective}
The central claim of the paper concerns quantum devices / superconducting circuits. Our objective was to reproduce the headline numerical/qualitative results within the constraints of the AI-driven Replication Project (24--72\,h, commodity GPU/CPU compute, no author contact). A replication plan is archived at \texttt{1983793-Simple-coplanar-waveguide-resonator-mask-targeting/replication\_plan.pdf}.

\section{Methodology}
We followed the standard Replication Project workflow: ingest the source PDF, extract method and data pointers, draft a replication plan, attempt the smallest end-to-end pipeline that produces a comparable number, then scale up where compute and data permit. Tools used in this replication are catalogued in the meta-paper's computational tools inventory (see \texttt{COMPUTATIONAL\_TOOLS\_INVENTORY.md}).

This paper went through the Tier-Lift pipeline; supplementary notes at \texttt{.openclaw/workspace/24h-progress/tier-lift-v2.5/1983793/REPORT.md}, \texttt{.openclaw/workspace/24h-progress/tier-lift-v2/1983793.md}.

\subsection*{Paper-directory README excerpt}
\small
\par\medskip\textbf{Simple coplanar waveguide resonator mask targeting metal-substrate interface}\par
\par$\bullet$\ **OSTI ID:** 1983793
\par$\bullet$\ **Rank:** \#7
\par$\bullet$\ **Replication Score:** 9/10

\par\medskip\textbf{\# Why This Paper}\par
This study uses finite-element electromagnetic simulations with open-source mask designs, structured outputs, and quantitative metrics, enabling full automation and validation by AI.

\par\medskip\textbf{\# Replication Plan}\par
AI would run finite-element simulations on the provided mask designs, vary geometric parameters, and calculate participation ratios and quality factors, comparing results to published figures and metrics.

\par\medskip\textbf{\# Status}\par
\par$\bullet$\ [ ] Paper reviewed
\par$\bullet$\ [ ] Data identified
\par$\bullet$\ [ ] Code implemented
\par$\bullet$\ [ ] Results validated
\normalsize

The replication was driven by an OpenClaw agent loop (Claude Opus / GPT-5 class models) issuing shell, Python, and (where applicable) GPU jobs. Compute usage and wall-time for individual papers are summarised in the meta-paper's resource table; not separately recorded for this paper unless noted in the Tier-Lift artefact. Sub-agents were spawned for replication-plan generation and (for papers in batches 1--4) for tier-lift extension runs.

\section{What was reproduced}
\begin{itemize}
\item Conformal-mapping CPW impedance formula
\item Kinetic-inductance correction (Z0 48.32 -\textgreater{} 49.26 ohm)
\item 8-resonator frequency plan (4.6, 5.0, 5.4, 5.8, 6.2, 6.6, 7.0, 7.4 GHz)
\item Quarter-wave resonator lengths
\item Participation ratio ordering p\_MS \textgreater{} p\_SA \textgreater{} p\_MA
\item Q\_c vs coupler length monotonic decrease over 3+ decades
\item Authors published GDS layout file (gdstk)
\item Aspect-ratio sweep s=3..12 um
\item Trench-depth sweep 0-300 nm
\item NEW: 3D FEM eigenmode for all 8 designed lengths in Palace
\item NEW: Quarter-wave harmonic series f1:f3:f5 = 1:3:5 to 0.1\% (correct BC physics)
\item NEW: Q-factor 1.087e5 = 1/(p\_sub*tan\_d) exactly across all 8 resonators
\item NEW: 8 frequencies match paper to \textless{}0.05\% after documented eps\_eff geometry correction
\end{itemize}

\section{What was missing or incomplete}
\begin{itemize}
\item Sub-nm interface meshing for direct FEM p\_MA/p\_MS/p\_SA recomputation
\item Full 8-resonator-on-feedline driven S-parameter sim with realistic coupler
\item Realistic isotropic-etch trench profile (curved)
\item Experimental Q\_i measurements and TLS loss-budget fitting
\item Kinetic-inductance two-fluid model in Palace eigenmode (vs purely dielectric loss)
\end{itemize}
The dominant blocker categories for this paper, mapped onto the meta-paper's 9-category friction taxonomy (\S4), are listed in \S\ref{sec:friction} below.

\section{Quantitative agreement}
Per-quantity numerical comparisons are not separately tabulated in the structured evaluation record for this paper beyond the agreement rationale below; where the rationale references specific numbers, they are reproduced verbatim. A full numerical comparison table is recorded in the per-replication artefacts (see \S\ref{sec:repro}). For papers below 8/8, the agreement rationale identifies the specific quantities that diverge.

\paragraph{Agreement rationale.} Quarter-wave harmonic series 1:3:5 reproduced to 0.1\% across 8 resonators (3D FEM diagnostic of correct short/open BC). Frequencies match paper design table to \textless{}0.05\% after eps\_eff geometry correction. Q-factor exactly matches analytical p\_sub * tan\_d prediction. Z0 with KI correction 49.26 ohm vs 50 ohm (0.7\%). p\_MS/p\_SA = 2.09 vs paper 2.00. Substrate participation 91.85\% vs 89\%. Aspect-ratio sweep gives ratio 1.7-2.6 across s=3..12 um. Trench-depth knee at d\textasciitilde{}50 nm consistent with paper choice of 100 nm.

\section{Score and friction-point tagging}\label{sec:friction}
\begin{itemize}
\item \textbf{Coverage:} 8/10 --- All v1 results retained (CPW Z0 via conformal mapping with KI correction, 8-res frequency plan, quarter-wave lengths, Q\_c vs coupler length, 2D FD participation ratios, gdstk inspection of authors GDS). NEW IN v2.5: full 3D FEM eigenmode simulation in Palace 0.13 (CPU, 4 MPI ranks) for all 8 designed resonator lengths covering 4.6-7.4 GHz. Quarter-wave standing-wave physics reproduced (f3/f1=3.000+-0.002, f5/f1=5.000+-0.005 across all 8 resonators). Q-factor 1.087e5 from Palace matches analytical Q = 1/(p\_sub * tan\_d) = 1.10e5 exactly. Geometry-induced 7.5\% eps\_eff offset (Palace 7.285 vs analytical 6.225 due to finite ground/air box) corrected post-hoc to \textless{}0.05\% per design freq. Still missing: sub-nm interface meshing for direct FEM p\_MA/p\_MS/p\_SA, full 8-resonator-on-feedline driven S-parameter sim, experimental Q\_i.
\item \textbf{Agreement:} 8/10 --- Quarter-wave harmonic series 1:3:5 reproduced to 0.1\% across 8 resonators (3D FEM diagnostic of correct short/open BC). Frequencies match paper design table to \textless{}0.05\% after eps\_eff geometry correction. Q-factor exactly matches analytical p\_sub * tan\_d prediction. Z0 with KI correction 49.26 ohm vs 50 ohm (0.7\%). p\_MS/p\_SA = 2.09 vs paper 2.00. Substrate participation 91.85\% vs 89\%. Aspect-ratio sweep gives ratio 1.7-2.6 across s=3..12 um. Trench-depth knee at d\textasciitilde{}50 nm consistent with paper choice of 100 nm.
\end{itemize}

\noindent\textbf{Friction points encountered (taxonomy tags):}
\begin{itemize}
\item None recorded.
\end{itemize}

\section{Follow-on questions}
\begin{itemize}
\item Not recorded.
\end{itemize}

\section{Reproducibility pointers}\label{sec:repro}
\begin{itemize}
\item \textbf{Paper directory:} \texttt{1983793-Simple-coplanar-waveguide-resonator-mask-targeting}
\item \textbf{Replication artefacts:}
\begin{itemize}
\item Replication artifacts in paper directory.
\end{itemize}
\item \textbf{Replication-dir hint from JSONL:} \texttt{\textasciitilde{}/Dropbox/REPLICATE-PROJECT/cpw-resonator/ + \textasciitilde{}/.openclaw/workspace/24h-progress/tier-lift-v2.5/1983793/}
\item \textbf{Status:} Not recorded
\end{itemize}

To re-run, change directory into the paper directory above and consult its \texttt{README.md} for the entry-point script. Most replications follow the convention \texttt{replication/run\_*.sh} or \texttt{replication/<subdir>/run.py}.

\end{document}
